% !TeX program = xelatex
\documentclass[a4paper,10pt]{article}

% ============================================================================
% សន្លឹកកិច្ចការសិស្ស ស-ខ្មៅ (Student Worksheet - Grade 8 3.2.5 SINT)
% បរិបទ៖ គ្មានស្លាយ LCD និងគ្មានម៉ាស៊ីនគិតលេខ — សិស្សបំពេញផ្ទាល់លើសន្លឹកកិច្ចការ ២ ទំព័រ A4
% ============================================================================
\usepackage[a4paper, left=0.85cm, right=0.85cm, top=0.70cm, bottom=0.95cm, footskip=0.48cm]{geometry}
\usepackage{amsmath, amssymb}
\usepackage{unicode-math}
\usepackage{fontspec}
\usepackage{array, colortbl}
\usepackage[table]{xcolor}
\usepackage{enumitem}
\usepackage{pifont}
\usepackage{fancyhdr}
\usepackage{lastpage}
\usepackage{tikz}
\usepackage[most]{tcolorbox}
\usetikzlibrary{calc, arrows.meta, decorations.pathreplacing, positioning, babel}

\makeatletter
\let\@expl@@@mark@update@singlecol@structures@@\relax
\makeatother

\providecommand{\NGSFindFonts}{%
  \IfFileExists{C:/Windows/Fonts/cambria.ttc}{%
    \def\NGSSysPath{C:/Windows/Fonts/}%
    \def\NGSUserPath{C:/Users/KIT/AppData/Local/Microsoft/Windows/Fonts/}%
  }{%
    \IfFileExists{./fonts/KhmerOSsiemreap.ttf}{\def\NGSSysPath{./fonts/}\def\NGSUserPath{./fonts/}}{%
    \IfFileExists{../fonts/KhmerOSsiemreap.ttf}{\def\NGSSysPath{../fonts/}\def\NGSUserPath{../fonts/}}{%
    \IfFileExists{../../fonts/KhmerOSsiemreap.ttf}{\def\NGSSysPath{../../fonts/}\def\NGSUserPath{../../fonts/}}{%
    \IfFileExists{../../../fonts/KhmerOSsiemreap.ttf}{\def\NGSSysPath{../../../fonts/}\def\NGSUserPath{../../../fonts/}}{%
    \IfFileExists{../../../../fonts/KhmerOSsiemreap.ttf}{\def\NGSSysPath{../../../../fonts/}\def\NGSUserPath{../../../../fonts/}}{%
    \IfFileExists{../../../../../fonts/KhmerOSsiemreap.ttf}{\def\NGSSysPath{../../../../../fonts/}\def\NGSUserPath{../../../../../fonts/}}{%
    \IfFileExists{../../../../../../fonts/KhmerOSsiemreap.ttf}{\def\NGSSysPath{../../../../../../fonts/}\def\NGSUserPath{../../../../../../fonts/}}{%
      \def\NGSSysPath{../../../../../../../fonts/}\def\NGSUserPath{../../../../../../../fonts/}%
    }}}}}}}%
  }%
}
\NGSFindFonts

\setmathfont[Path={\NGSSysPath}, FontIndex=1, Scale=1.02]{cambria.ttc}
\defaultfontfeatures{Ligatures=TeX}

\setmainfont[
  Path={\NGSUserPath}, Extension=.ttf, Script=Khmer,
  BoldFont={KhmerOSsiemreap}, BoldFeatures={FakeBold=1.8},
  ItalicFont={KhmerOSsiemreap}, ItalicFeatures={FakeSlant=0.22}
]{KhmerOSsiemreap}

\newfontfamily\muol[Path={\NGSUserPath}, Extension=.ttf, Script=Khmer]{KhmerOSmuollight}
\newfontfamily\TimesEn[Path={\NGSSysPath}, Extension=.ttf, BoldFont={timesbd}, ItalicFont={timesi}]{times}

\XeTeXlinebreaklocale "km"
\XeTeXlinebreakskip = 0pt plus 1pt minus 0.1pt

\definecolor{panelgray}{gray}{0.93}
\definecolor{headgray}{gray}{0.85}

\pagestyle{fancy}
\fancyhf{}
\renewcommand{\headrulewidth}{0pt}
\fancyfoot[L]{\fontsize{7.5pt}{9.5pt}\selectfont អនុវិទ្យាល័យព្រះក្រសាំង \en{|} គណិតវិទ្យាថ្នាក់ទី៨ \en{|} មេរៀនទី៣៖ ២.៥ ការប្រាក់សាមញ្ញ (គ្រូ៖ ហុក ប៊ុនថែន)}
\fancyfoot[R]{\fontsize{7.8pt}{9.5pt}\selectfont ទំព័រ \en{\thepage/\pageref{LastPage}}}

\setlength{\parindent}{0pt}
\setlength{\parskip}{0pt}
\setlength{\tabcolsep}{3.8pt}

\newcommand{\en}[1]{{\TimesEn #1}}
\newcommand{\enbf}[1]{{\TimesEn\bfseries #1}}
\newcommand{\unit}[1]{\,\text{\en{#1}}}
\newcommand{\cnum}[1]{\mbox{\fontsize{10pt}{12pt}\selectfont\ding{\numexpr181+#1\relax}}}
\newcommand{\chkbox}{\tikz[baseline=-0.12ex]\draw[black, line width=0.8pt] (0,0) rectangle (0.24cm,0.24cm);}
\newcommand{\blank}[1][1.5cm]{\makebox[#1]{\dotfill}}
\newcommand{\ansbox}[1][0.85cm]{\,\tikz[baseline=-0.5ex]\draw[black, line width=0.65pt, fill=white] (0,0) rectangle (#1, 0.42cm);\,}
\newcommand{\secbar}[2]{%
  \begin{tcolorbox}[enhanced, sharp corners, colback=headgray, colframe=black, boxrule=0.9pt, boxsep=1pt, top=2pt, bottom=2pt, left=5pt, right=5pt, before skip=4pt, after skip=3pt]
    {\muol\fontsize{9.2pt}{13pt}\selectfont #1}\hfill{\fontsize{8pt}{11pt}\selectfont\textbf{#2}}
  \end{tcolorbox}%
}

\newcolumntype{L}[1]{>{\raggedright\arraybackslash}p{#1}}
\newcolumntype{C}[1]{>{\centering\arraybackslash}p{#1}}

\begin{document}
\fontsize{9.6pt}{14.4pt}\selectfont

% ============================================================================
% ទំព័រទី ១៖ E1 + E2 (POE 1 & POE 2) + E3 សំយោគរូបមន្ត និងធនាគារហេតុផល
% ============================================================================
\begin{tcolorbox}[enhanced, sharp corners, colback=white, colframe=black, boxrule=1.0pt, boxsep=2pt, top=3pt, bottom=3pt, left=6pt, right=6pt]
  \begin{center}
    {\muol\fontsize{10.0pt}{15pt}\selectfont អនុវិទ្យាល័យព្រះក្រសាំង \enbf{--} សន្លឹកកិច្ចការសិស្ស \enbf{(}មេរៀនទី៣៖ ២\enbf{.}៥ ការប្រាក់សាមញ្ញ\enbf{)}}\\[3pt]
    {\fontsize{8.8pt}{12pt}\selectfont
    ឈ្មោះសិស្ស / ដៃគូ៖ \blank[5.2cm] \hspace{0.25cm}
    ថ្នាក់ទី៨៖ \blank[1.1cm] \hspace{0.25cm}
    ក្រុមទី៖ \blank[1.0cm] \hspace{0.25cm}
    កាលបរិច្ឆេទ៖ ០១/១២/២០២៦}
  \end{center}
\end{tcolorbox}

% --- E1 + E2 ---
\secbar{ផ្នែកទី១៖ បេសកកម្មទស្សន៍ទាយ \enbf{(E1)} និងការរុករកជាដៃគូ \enbf{(E2 -- POE 1 \& POE 2)}}{រយៈពេល ១៤ នាទី}

\begin{tcolorbox}[enhanced, sharp corners, colback=panelgray, colframe=black, boxrule=0.75pt, boxsep=1.5pt, top=2.5pt, bottom=2.5pt, left=5pt, right=5pt]
  \textbf{បេសកកម្ម \en{E1} (គូសធីក \chkbox\ មុនពិសោធន៍)៖}
  (១) ដាក់ប្រាក់ដើម $2\,000\,000\text{៛}$ អត្រា $15\%$/ឆ្នាំ រយៈពេល $3$ ឆ្នាំ — ការប្រាក់ស្មើ៖ \chkbox\ $300\,000\text{៛}$ \quad \chkbox\ $900\,000\text{៛}$\\
  (២) ខ្ចីប្រាក់ដើម $6\,000\,000\text{៛}$ អត្រា $12\%$/ឆ្នាំ រយៈពេល $8$ ខែ — ការប្រាក់ស្មើ៖ \chkbox\ $5\,760\,000\text{៛}$ \quad \chkbox\ $480\,000\text{៛}$
\end{tcolorbox}

\vspace{3pt}
\textbf{ក. សកម្មភាពទី ១ (\en{POE 1})៖ កសាងរូបមន្តការប្រាក់សាមញ្ញ $I = \frac{PRT}{100}$ ពីកំណើនថេរប្រចាំឆ្នាំ (ទំ.៣៥)}

{\renewcommand{\arraystretch}{1.46}
\begin{tabular}{|C{2.40cm}|C{3.30cm}|C{2.70cm}|L{5.50cm}|C{3.70cm}|}
  \hline
  \rowcolor{headgray}
  \textbf{រយៈពេល $T$} & \textbf{ប្រាក់ដើម $P$} & \textbf{អត្រា $R$ (/ឆ្នាំ)} & \textbf{ការគណនាការប្រាក់សាមញ្ញ $I$} & \textbf{ប្រាក់សរុប $A = P + I$} \\
  \hline
  $1$ ឆ្នាំ & $2\,000\,000\text{៛}$ & $15\%$ & $2\,000\,000 \times \dfrac{15}{100} \times 1 = \blank[1.5cm]\text{៛}$ & $2\,300\,000\text{៛}$ \\
  \hline
  $2$ ឆ្នាំ & $2\,000\,000\text{៛}$ & $15\%$ & $2\,000\,000 \times \dfrac{15}{100} \times 2 = \blank[1.5cm]\text{៛}$ & $\blank[2.0cm]\text{៛}$ \\
  \hline
  $3$ ឆ្នាំ & $2\,000\,000\text{៛}$ & $15\%$ & $\dfrac{2\,000\,000 \times 15 \times \blank[0.6cm]}{100} = \blank[1.5cm]\text{៛}$ & $\blank[2.0cm]\text{៛}$ \\
  \hline
\end{tabular}}

\vspace{3pt}
$\bullet$ \textbf{ពន្យល់ (\en{E})៖} សម្រាប់រយៈពេល $T$ ឆ្នាំ ការប្រាក់សាមញ្ញគណនាតាមរូបមន្ត $I = \dfrac{P \times R \times \blank[0.8cm]}{100}$ ។

\vspace{4pt}
\textbf{ខ. សកម្មភាពទី ២ (\en{POE 2})៖ ច្រកប្ដូររយៈពេលខែជាប្រភាគនៃឆ្នាំ ($T = \frac{m}{12}$) និងការទាញរកអញ្ញាតបញ្ច្រាស}

{\renewcommand{\arraystretch}{1.46}
\begin{tabular}{|L{4.50cm}|C{3.50cm}|L{6.20cm}|C{3.40cm}|}
  \hline
  \rowcolor{headgray}
  \textbf{ករណីសិក្សា} & \textbf{ប្ដូររយៈពេលជាឆ្នាំ ($T$)} & \textbf{ការជំនួសចូលរូបមន្តគណនាដោយដៃ} & \textbf{លទ្ធផលត្រឹមត្រូវ} \\
  \hline
  ១) ខ្ចី $6\,000\,000\text{៛}$, $R=12\%$, $8$ ខែ & $8\text{ ខែ} = \dfrac{8}{12} = \blank[0.8cm]$ ឆ្នាំ & $I = \dfrac{6\,000\,000\times 12\times \blank[0.8cm]}{100} = \blank[1.3cm]$ & $\blank[1.6cm]\text{៛}$ \\
  \hline
  ២) ដាក់ $5\,000\,000\text{៛}$, $2$ ឆ្នាំ បាន $I=800\,000\text{៛}$ & $T = 2$ ឆ្នាំ (ស្គាល់) & $R = \dfrac{100 \times I}{P \times T} = \dfrac{100 \times 800\,000}{5\,000\,000 \times 2}$ & $R = \blank[1.0cm]\%$ \\
  \hline
\end{tabular}}

\vspace{3pt}
$\bullet$ \textbf{ពន្យល់ (\en{E})៖} ហេតុអ្វីជំនួស $T=8$ សម្រាប់ $8$ ខែខុស? ព្រោះអត្រា $R$ គិតក្នុង $1$ ឆ្នាំ ដូច្នេះ $T=8$ មានន័យថា \blank[1.8cm] មិនមែន $8$ ខែទេ!

% --- E3 ---
\secbar{ផ្នែកទី២៖ សំយោគរូបមន្តគោល \enbf{(E3)} និងធនាគារហេតុផល \enbf{B\&W} \cnum{1}\enbf{--}\cnum{5}}{សៀវភៅពុម្ព ទំ.៣៥–៣៦}

\begin{tabular}{|L{9.20cm}|L{9.20cm}|}
  \hline
  \textbf{១. រូបមន្តការប្រាក់សាមញ្ញ និងរូបមន្តបញ្ច្រាស៖}\par\vspace{2pt}
  $\bullet$ \textbf{ការប្រាក់សាមញ្ញ៖} $I = \dfrac{P \cdot R \cdot \blank[0.7cm]}{100}$ ($T$ គិតជា\textbf{ឆ្នាំ})\par\vspace{2.5pt}
  $\bullet$ \textbf{ប្រាក់សរុបចុងគ្រា៖} $A = P + \blank[0.8cm]$\par\vspace{2.5pt}
  $\bullet$ \textbf{ទាញរកអត្រា $R$ ($\%$ / ឆ្នាំ)៖} $R = \dfrac{100 \cdot I}{\blank[1.5cm]}$\par\vspace{2.5pt}
  $\bullet$ \textbf{ទាញរកប្រាក់ដើម ឬពេល៖} $P = \dfrac{100 I}{R T}$ \;;\; $T = \dfrac{100 I}{\blank[1.2cm]}$
  &
  \textbf{២. ធនាគារហេតុផល \en{B\&W} (\cnum{1}--\cnum{5})៖}\par\vspace{2pt}
  \cnum{1} រូបមន្តគោល៖ $I = \frac{PRT}{100}$ គិតការប្រាក់ថេរលើប្រាក់ដើម $P$ ។\par\vspace{2pt}
  \cnum{2} ច្រកប្ដូរខែជាឆ្នាំ៖ បើរយៈពេល $m$ ខែ ត្រូវប្ដូរ $T = \frac{m}{12}$ ឆ្នាំសិន ។\par\vspace{2pt}
  \cnum{3} ប្រាក់សរុបចុងគ្រា៖ ប្រាក់ដើមបូកការប្រាក់ $A = P + I$ ។\par\vspace{2pt}
  \cnum{4} ទាញរកអត្រាការប្រាក់៖ $R = \frac{100 I}{P T}$ (គិតជា $\%$ ក្នុងមួយឆ្នាំ) ។\par\vspace{2pt}
  \cnum{5} ទាញរកដើម ឬពេល៖ $P = \frac{100 I}{R T}$ និង $T = \frac{100 I}{P R}$ ។ \\
  \hline
\end{tabular}

\vspace{3pt}
\begin{tcolorbox}[enhanced, sharp corners, colback=panelgray, colframe=black, boxrule=0.65pt, boxsep=1.5pt, top=2pt, bottom=2pt, left=5pt, right=5pt]
  \textbf{ប្រយ័ត្នអន្ទាក់៖} មុនជំនួសលេខចូល $I = \frac{PRT}{100}$ ត្រូវពិនិត្យឯកតារយៈពេល $T$ ជានិច្ច ($3\text{ ខែ}=\frac{1}{4}$ ឆ្នាំ, $6\text{ ខែ}=\frac{1}{2}$ ឆ្នាំ, $8\text{ ខែ}=\frac{2}{3}$ ឆ្នាំ, $9\text{ ខែ}=\frac{3}{4}$ ឆ្នាំ)!
\end{tcolorbox}

\vspace{3pt}
\textbf{៣. លំហាត់ត្រួតពិនិត្យរហ័ស៖ ប្ដូរខែជាឆ្នាំ និងគណនាការប្រាក់សាមញ្ញដោយសម្រួលសូន្យ}

\vspace{2pt}
{\renewcommand{\arraystretch}{1.60}
\begin{tabular}{|C{4.40cm}|C{4.40cm}|C{4.40cm}|C{4.40cm}|}
  \hline
  \rowcolor{panelgray}
  $6\text{ ខែ} = \blank[0.8cm]$ ឆ្នាំ & $4\text{ ខែ} = \blank[0.8cm]$ ឆ្នាំ & $9\text{ ខែ} = \blank[0.8cm]$ ឆ្នាំ & $18\text{ ខែ} = \blank[0.8cm]$ ឆ្នាំ \\
  \hline
  $P=1\,000\$, R=5\%, 2\text{ឆ្នាំ}\to I=\blank[0.6cm]$ & $P=2\,000\$, R=10\%, 3\text{ឆ្នាំ}\to I=\blank[0.6cm]$ & $P=400\$, R=8\%, \frac{1}{2}\text{ឆ្នាំ}\to I=\blank[0.6cm]$ & $P=500\$, I=50\$, 2\text{ឆ្នាំ}\to R=\blank[0.5cm]\%$ \\
  \hline
  $P=3\,000\$, R=6\%, 4\text{ខែ}\to I=\blank[0.6cm]$ & $P=800\$, R=10\%, 9\text{ខែ}\to I=\blank[0.6cm]$ & $I=120\$, R=6\%, 2\text{ឆ្នាំ}\to P=\blank[0.6cm]$ & $P=1\,000\$, R=12\%, I=360\$\to T=\blank[0.5cm]$ \\
  \hline
  $P=50\,000\text{៛}, R=10\%, 1\text{ឆ្នាំ}\to A=\blank[0.6cm]$ & $P=600\$, R=5\%, 4\text{ឆ្នាំ}\to A=\blank[0.6cm]$ & $P=1\,200\$, R=10\%, 6\text{ខែ}\to I=\blank[0.6cm]$ & $I=90\$, P=300\$, 3\text{ឆ្នាំ}\to R=\blank[0.5cm]\%$ \\
  \hline
\end{tabular}}

\newpage

% ============================================================================
% ទំព័រទី ២៖ E4 លំហាត់គំរូ + E5 ប្រតិបត្តិ + វិនិច្ឆ័យ A-D + Exit Check + កិច្ចការផ្ទះ
% ============================================================================
\secbar{ផ្នែកទី៣៖ ពង្រីកចំណេះដឹងលើលំហាត់គំរូប្រាក់បញ្ញើមីងផុន \enbf{(E4 --} ទំ\enbf{.}៣៦\enbf{)}}{រយៈពេល ១០ នាទី}

\textbf{លំហាត់គំរូ (សៀវភៅពុម្ព ទំ.៣៦ — មីងផុនដាក់ប្រាក់នៅធនាគារ)៖} មីងផុនយកប្រាក់ $\mathbf{15\,000\,000\text{៛}}$ ទៅដាក់ធនាគារដោយគិតការប្រាក់សាមញ្ញ៖ (ក) បើអត្រាការប្រាក់ $\mathbf{6\%}$ ក្នុងមួយឆ្នាំ រយៈពេល $\mathbf{5}$ ឆ្នាំ រកប្រាក់ការសរុប និងប្រាក់សរុបដែលគាត់ទទួលបាន ។ (ខ) បើគាត់ដាក់រយៈពេល $\mathbf{2}$ ឆ្នាំ ហើយទទួលបានប្រាក់ការ $\mathbf{3\,000\,000\text{៛}}$ តើធនាគារគិតអត្រាការប្រាក់ប៉ុន្មានភាគរយក្នុងមួយឆ្នាំ?

{\renewcommand{\arraystretch}{1.50}
\begin{tabular}{|L{4.50cm}|L{9.10cm}|C{4.20cm}|}
  \hline
  \rowcolor{headgray}
  \textbf{សំណួរ និងទិន្នន័យ} & \textbf{ការជំនួសរូបមន្ត និងគណនាដោយដៃ} & \textbf{លទ្ធផល និងហេតុផល} \\
  \hline
  \textbf{ក.១ រកប្រាក់ការ ($I$)} ($P=15\text{លាន}, R=6\%, T=5$) & $I = \dfrac{P R T}{100} = \dfrac{15\,000\,000 \times 6 \times 5}{100} = \blank[2.0cm]\text{៛}$ & ប្រាក់ការ $5$ ឆ្នាំ (\cnum{1}) \\
  \hline
  \textbf{ក.២ រកប្រាក់សរុប ($A$)} \textit{(កែតម្រូវលេខបោះពុម្ពទំ.៣៦)} & $A = P + I = 15\,000\,000 + \blank[1.8cm] = \blank[2.0cm]\text{៛}$ & ប្រាក់សរុបចុងគ្រា (\cnum{3}) \\
  \hline
  \textbf{ខ. រកអត្រាការប្រាក់ ($R$)} ($T=2\text{ ឆ្នាំ}, I=3\,000\,000\text{៛}$) & $R = \dfrac{100 \times I}{P \times T} = \dfrac{100 \times 3\,000\,000}{15\,000\,000 \times \blank[0.7cm]} = \blank[1.2cm]\%$ & អត្រា $\blank[0.9cm]\%$/ឆ្នាំ (\cnum{4}) \\
  \hline
\end{tabular}}

\vspace{4pt}
$\bullet$ \textbf{សំណួរឆ្លុះបញ្ចាំង \en{(EIA Looking-Back)}៖} ហេតុអ្វីក្នុងការប្រាក់សាមញ្ញ ការប្រាក់ក្នុងឆ្នាំទី ១ ឆ្នាំទី ២ និងឆ្នាំទី ៥ ស្មើ $900\,000\text{៛}$ ដូចគ្នារាល់ឆ្នាំ? \blank[2.4cm]

% --- E5 ---
\secbar{ផ្នែកទី៤៖ វាយតម្លៃលទ្ធផលសិក្សា \enbf{(E5 --} ប្រតិបត្តិ ទំ\enbf{.}៣៧ \enbf{+} វិនិច្ឆ័យ \enbf{A--D +} សំបុត្រចេញ\enbf{)}}{រយៈពេល ១០ នាទី}

\textbf{ក. លំហាត់ប្រតិបត្តិ (សៀវភៅពុម្ព ទំ.៣៧ និងពង្រឹងខែ-ឆ្នាំ — ធ្វើជាដៃគូ រួចប្ដូរសន្លឹកផ្ទៀង \en{Pause-Review})៖}

{\renewcommand{\arraystretch}{1.50}
\begin{tabular}{|L{8.90cm}|L{8.90cm}|}
  \hline
  \textbf{១) ប្រតិបត្តិ (ទំ.៣៧ — ប្រាក់បញ្ញើរបស់ជីកទ្រី)៖}\par\vspace{2pt}
  ជីកទ្រីយកប្រាក់ $\mathbf{24\,000\,000\text{៛}}$ ទៅដាក់ធនាគារដោយគិតការប្រាក់សាមញ្ញតាមអត្រា $\mathbf{16\%}$ ក្នុងមួយឆ្នាំ រយៈពេល $\mathbf{4}$ ឆ្នាំ ។ គណនាប្រាក់ការ និងប្រាក់សរុប៖\par\vspace{2pt}
  $\bullet$ $I = \dfrac{24\,000\,000 \times 16 \times 4}{100} = \blank[2.0cm]\text{៛}$ (\cnum{1})\par\vspace{2pt}
  $\bullet$ ប្រាក់សរុប $A = 24\,000\,000 + \blank[1.6cm] = \blank[1.9cm]\text{៛}$ (\cnum{3})
  &
  \textbf{២) ពង្រឹង (ការប្រាក់សាមញ្ញរយៈពេលគិតជាខែ)៖}\par\vspace{2pt}
  កសិករម្នាក់ខ្ចីប្រាក់ $\mathbf{9\,000\,000\text{៛}}$ ពីគ្រឹះស្ថានមីក្រូហិរញ្ញវត្ថុ តាមអត្រាការប្រាក់សាមញ្ញ $\mathbf{10\%}$ ក្នុងមួយឆ្នាំ រយៈពេល $\mathbf{8}$ \textbf{ខែ} ។ គណនាប្រាក់ការ និងប្រាក់ត្រូវសងសរុប៖\par\vspace{2pt}
  $\bullet$ ប្ដូរខែជាឆ្នាំ៖ $T = \dfrac{8}{12} = \blank[0.9cm]$ ឆ្នាំ (\cnum{2})\par\vspace{2pt}
  $\bullet$ $I = \dfrac{9\,000\,000\times 10\times \blank[0.7cm]}{100} = \blank[1.4cm]\text{៛} \Rightarrow A = \blank[1.6cm]\text{៛}$ \\
  \hline
\end{tabular}}

\vspace{5pt}
\textbf{ខ. គូសវង់ជុំវិញជម្រើសត្រឹមត្រូវតែ ១ គត់ ក្នុងចំណោម \en{A, B, C, D} រួចសរសេរលេខហេតុផល \cnum{1}--\cnum{5} នៃជម្រើសខុស៖}

{\renewcommand{\arraystretch}{1.36}
\begin{tabular}{|C{4.40cm}|C{4.40cm}|C{4.40cm}|C{4.40cm}|}
  \hline
  \cellcolor{panelgray}\textbf{ជម្រើស \en{A}} & \cellcolor{panelgray}\textbf{ជម្រើស \en{B}} & \cellcolor{panelgray}\textbf{ជម្រើស \en{C}} & \cellcolor{panelgray}\textbf{ជម្រើស \en{D}} \\
  \hline
  $P=9\text{លាន}, 10\%, 8\text{ ខែ}$ $\Rightarrow I = 600\,000\text{៛}$ & $P=9\text{លាន}, 10\%, T=8$ $\Rightarrow I = 7\,200\,000\text{៛}$ & ជីកទ្រីដាក់ $24\text{លាន}$ សរុប $A = 15\,360\,000\text{៛}$ & រកអត្រាការប្រាក់តាម $R = \dfrac{P T}{100 I}$ \\[3pt]
  \chkbox\ ត្រូវ \quad \chkbox\ ខុស (\blank[0.7cm]) & \chkbox\ ត្រូវ \quad \chkbox\ ខុស (\blank[0.7cm]) & \chkbox\ ត្រូវ \quad \chkbox\ ខុស (\blank[0.7cm]) & \chkbox\ ត្រូវ \quad \chkbox\ ខុស (\blank[0.7cm]) \\
  \hline
\end{tabular}}

\vspace{5pt}
\begin{tcolorbox}[enhanced, sharp corners, colback=white, colframe=black, boxrule=0.85pt, boxsep=2.5pt, top=3.5pt, bottom=3.5pt, left=5pt, right=5pt]
  {\muol\fontsize{8.8pt}{12.5pt}\selectfont សំបុត្រចេញរហ័ស \enbf{(Exit Check --} ៦០ វិនាទី\enbf{)}៖}\\[2pt]
  (១) $P=5\,000\,000\text{៛}, R=12\%$/ឆ្នាំ រយៈពេល $6$ ខែ៖ $I = \blank[1.8cm]\text{៛}$ \quad\en{|}\quad
  (២) $P=4\,000\,000\text{៛}, T=2$ ឆ្នាំ បាន $I=640\,000\text{៛}$៖ $R = \blank[1.1cm]\%$
\end{tcolorbox}

\vspace{3pt}
\begin{tcolorbox}[enhanced, sharp corners, colback=panelgray, colframe=black, boxrule=0.75pt, boxsep=2pt, top=3pt, bottom=3pt, left=5pt, right=5pt]
  {\muol\fontsize{8.8pt}{12.5pt}\selectfont កិច្ចការផ្ទះ \enbf{(}សៀវភៅគោល ទំព័រ ៤២ លំហាត់ \enbf{\#16} ក\enbf{,} ខ\enbf{,} គ\enbf{)}៖}\\[2pt]
  $\bullet$ \textbf{លំហាត់ \#16 (ក, ខ, គ — ទំ.៤២)៖} គណនាការប្រាក់សាមញ្ញ ប្រាក់សរុប អត្រាការប្រាក់ និងប្ដូររយៈពេលខែជាឆ្នាំឱ្យបានត្រឹមត្រូវ ។\\
  $\bullet$ \textbf{ត្រៀមមេរៀនបន្ទាប់៖} អានមុន \textbf{២.៦. ការប្រាក់សមាស (ទំ.៣៧–៣៨)} ។
\end{tcolorbox}

\end{document}
